% SPDX-License-Identifier: GPL-3.0-or-later
% Copyright (C) 2026 Alessandro Greco
\chapter{torspec 协议规范阅读指南}\label[appendice]{cap:C-specifiche}
\aggiornatoal{2026 年 9 月 23 日}{torspec 提交 \texttt{8e1f245}}

Tor Project 的规范是本书几乎每一章的主要资料来源。它们内容全面，却不是按照教科书形式编写的。二十多年来，规范最初以互相独立的文本文件维护，直到较近时期才汇总到同一个网站。规范自己的约定文档也承认，其一致性仍在持续完善中。本附录为希望读完本书后继续研究规范的读者提供一份导航图。

% ---------------------------------------------------------------------------
\section{从哪里找到规范？}\label{sec:C-dove}

规范发布在 \url{https://spec.torproject.org/}，并与提案一起在 Tor Project 的 \texttt{torspec} GitLab 仓库维护。规范是 \file{spec/} 目录中的 Markdown 文件：较小的主题由单独文件描述，内容较多的文档则使用目录。文件 \file{permalinks.md} 列出了文档的稳定地址，例如 \url{https://spec.torproject.org/tor-spec}；引用规范时应优先使用这些链接。本书一直采用 \spec{document}{section} 形式引用规范文档及章节，保留网站上使用的原始章节标题；提案则写作\enquote{提案~NNN}。

% ---------------------------------------------------------------------------
\section{约定}\label{sec:C-convenzioni}

规范的引言约定文档规定了一些通用规则，除非具体章节另有说明：
\begin{itemize}
  \item \texttt{MUST}、\texttt{SHOULD}、\texttt{MAY} 及其否定形式，采用 RFC~2119 中的含义，分别表示强制要求、强烈建议和允许的选项。
  \item 长度以八比特字节为单位，\enquote{kibibyte} 等二进制前缀代表 1024 的倍数。
  \item 多字节整数采用大端序，即\enquote{网络字节序}：十进制整数 4660 以两个字节表示为 \texttt{12 34}。
\end{itemize}
本书前面遇到过两个重要例外：Curve25519 与 Ed25519 的密钥及点采用小端序（\cref{cap:B-crittografia}）；部分目录格式则是包含以 base64 编码的二进制数据的可读文本。

% ---------------------------------------------------------------------------
\section{规范文档地图}\label{sec:C-mappa}

\begin{center}\small
\begin{tabularx}{\textwidth}{@{}l>{\raggedright\arraybackslash}Xl@{}}
\toprule
文档 & 内容 & 相关章节 \\
\midrule
\texttt{tor-spec} & 核心协议：通道、信元、电路、握手、中继加密、流、流量控制和拥塞控制 & \ref{cap:07-canali}--\ref{cap:11-stream} \\
\texttt{cert-spec} & Tor 的 Ed25519 证书 & \ref{cap:07-canali} \\
\texttt{dir-spec} & 目录权威、投票、共识、描述符、微描述符与 \texttt{extra-info} 文档 & \ref{cap:12-directory}, \ref{cap:24-metrics} \\
\texttt{path-spec} & 路径选择及其约束 & \ref{cap:13-percorso} \\
\texttt{guard-spec} & 守卫选择算法 & \ref{cap:13-percorso} \\
\texttt{vanguards-spec} & 完整与轻量级 Vanguard 保护机制 & \ref{cap:13-percorso} \\
\texttt{rend-spec} & v3 洋葱服务：密钥、描述符、引介及汇合 & \ref{cap:14-onion-services}, \ref{cap:21-onion-pratica} \\
\texttt{hspow-spec} & 洋葱服务工作量证明 & \ref{cap:15-dos} \\
\texttt{dos-spec} & 拒绝服务攻击防御 & \ref{cap:15-dos} \\
\texttt{pt-spec} & tor 与可插拔传输程序之间的接口 & \ref{cap:16-censura} \\
\texttt{padding-spec} & 通道与电路填充 & \ref{cap:07-canali}, \ref{cap:26-attacchi} \\
\texttt{control-spec} & ControlPort 控制协议 & \ref{cap:18-stem} \\
\texttt{socks-extensions} & Tor 专用的 SOCKS 扩展 & \ref{cap:11-stream}, \ref{cap:28-opsec} \\
\texttt{param-spec} & 共识网络参数及其默认值 & \ref{cap:11-stream}, \ref{cap:15-dos} \\
\texttt{address-spec} & 特殊地址名称，例如 \texttt{.onion} 地址 & \ref{cap:14-onion-services} \\
\texttt{srv-spec} & 目录权威节点的共享随机值 & \ref{cap:14-onion-services} \\
\bottomrule
\end{tabularx}
\end{center}

% ---------------------------------------------------------------------------
\section{如何阅读规范章节？}\label{sec:C-lettura}

以下建议来自本书的编写过程：
\begin{enumerate}
  \item \textbf{先了解格式。}描述二进制消息的章节会以表格列出字段名称及长度。阅读规则之前，应像本书的格式示意图那样，先在纸上重建消息格式。
  \item \textbf{区分当前功能与过时内容。}许多文档仍保留 TAP、链路协议版本 1--3 和 v2 洋葱服务等历史功能。过时资料通常集中在 \file{obsolete-channels.md}、\file{historical/} 等章节或文件，但并非总是如此。存在疑问时，应检查源码对应分支是否仍然存在。
  \item \textbf{与实现相比较。}规范表达设计意图，tor 和 Arti 的源码揭示实际行为。本书研究期间，一些差异很有启发：单跳电路中未出现的通道填充（\cref{cap:07-canali}），以及出口中继返回的流关闭原因（\cref{sec:lab-relay}）。
  \item \textbf{检查提案状态。}有些提案在纳入正式规范前就已经实现，另一些仍在讨论中，尚无实现（\cref{sec:fut-processo}）。
  \item \textbf{进行实验验证。}本书的私有实验网络可用于检查某个字段是否按文档描述运行。这也是全部实验共同采用的方法。
\end{enumerate}

% ---------------------------------------------------------------------------
\section{Arti 与规范}\label{sec:C-arti}

Arti 的源码经常在注释和文档中引用规范，并指出其中模糊的要求或与 tor 的差异。它支持的子协议表格（\cref{sec:arti-supporto}）说明两种实现共同支持哪些功能，因此也显示了哪些规范内容描述的是仍在实际运行的\enquote{活协议}。
